% Authorship: Claude-drafted, author-revised
\section{Introduction}

Where do the determinants of legal content come from? The question sounds
technical, yet it sits at the center of a half-century jurisprudential
impasse.
On H.L.A.\ Hart's account, law's content is determined by social facts, by
convergent official practice, rather than by moral
truths.\footnote{H.L.A.\ Hart, \textsc{The Concept of Law} 100--10 (3d
ed.\ 2012).} On its face the thesis is attractive. On that account, law can exist in
morally deficient regimes, and legal disputes can be settled without first
resolving deep moral disagreements.

On the details, though, Hart's framework is conspicuously silent. How much
convergence is ``enough'' for a legal norm to exist? Where judicial practice
pulls in conflicting directions, with some judges treating a constitutional
value as dispositive and others weighing it against competing principles,
does the system generate contradictory law, or no law at all? Hart's answer
was candid and unsatisfying. In such ``hard cases,'' he said, the law simply
runs out.\footnote{\textit{Id.}\ at 124--36 (discussing the ``open texture''
of law and the inevitability of judicial discretion in hard cases).} In those cases, as Hart saw it, judges exercise discretion, and their
decisions make new law rather than discovering law that was already there.\footnote{\textit{Id.}\ at 272
(Postscript).}

To Ronald Dworkin, that was a fatal concession. If law runs out in precisely
the cases that matter most, where rights hang in the balance and the parties
most need resolution, then Hart's positivism delivers too little law to do
the work law must do.\footnote{Ronald Dworkin, \textsc{Law's Empire} 37--43
(1986).} Against that, Dworkin offered law as integrity: on any question the law is
whatever follows from the best moral interpretation of the community's
institutional history. Hard cases \textit{do} have right answers on that
view. Finding them, however, requires constructive moral reasoning of a kind
that resists formal specification.\footnote{\textit{Id.}\ at
225--75.}

For half a century the result has been an impasse, and a shabby place for
jurisprudence to have settled. Hart's theory is formal enough to model but appears to undergenerate law.
Dworkin's generates law enough but resists the formalization that would let
anyone test it systematically.\footnote{For the canonical framing, see Scott J.\ Shapiro,
\textsc{Legality} 25--52 (2011). For more recent assessments, see Jules L.\
Coleman, \textsc{The Practice of Principle} 65--72 (2001); Mark Greenberg,
\textit{The Moral Impact Theory of Law}, 123 \textsc{Yale L.J.}\ 1288,
1290--95 (2014).}

Mitchell Berman's principled positivism offers a way
forward.\footnote{Mitchell N.\ Berman, \textit{How Practices Make
Principles and How Principles Make Rules}, 28 \textsc{J.\ Ethics \& Soc.\
Phil.}\ 299 (2024) [hereinafter Berman, \textit{Practices}]. Berman is
developing these ideas further in a book in progress. See Mitchell N.\
Berman, \textsc{To Say What the Law Is: The Constitution, the Supreme
Court, and the Nature of Law} (forthcoming) (manuscript on file with
author). I am grateful to Professor Berman for extensive discussions
clarifying the theory and responding to the computational findings
reported here.} Like Hart, Berman identifies law through social facts and
convergent official practice. Unlike Hart, he does not rely on consensus.
Instead, legal principles \textit{aggregate} through Simple Additive
Weighting. Each principle contributes legal force in proportion to two
quantities: its weight, meaning how much it counts in the legal system
generally, and its activation, meaning how strongly it bears on the case at
hand. In a given
case the net legal force is:
%
\begin{equation}
F = \sum_{i=1}^{n} w_i \times a_i
\end{equation}
%
A case is determinate when the principles together generate net force in one
direction. How much force is required, and indeed whether the theory needs
any threshold beyond the bare balance, is among the questions taken up
below.\footnote{Berman, \textit{Practices}, \textit{supra} note 7, at
328--31.}

The proposal has three virtues. Because principles need not achieve
consensus, it offers more law than Hart: several principles, each carrying
only partial support, can aggregate to a determinate outcome where no single
principle commands universal assent. Because the aggregation is mechanical,
a matter of multiplying weights by activations and summing, it avoids
Dworkin's demand for moral interpretation. And it explains something neither Hart nor Dworkin captures well. In
practice, judges appear to treat conflicting considerations as carrying
\textit{different weights} to be \textit{balanced}, rather than as
generating contradictions or as demanding interpretive resolution.

Berman develops these claims philosophically rather than quantitatively,
and the theory leaves four questions open. Each bears on the proposal's
success.

First, how precise must the weights be? Berman argues that rough weights
suffice.\footnote{Berman, \textit{Practices}, \textit{supra} note 7, at
328--30 (responding to Stephen Perry's objection that judicial practice
cannot deliver precise enough weights). See Stephen R.\ Perry, \textit{Two Models of Legal Principles}, 82
\textsc{Iowa L.\ Rev.}\ 787, 794 (1997) (doubting that custom ``could be
sufficiently nuanced as to be able to assign determinate weights to
individual principles''), which is the passage Berman answers.} He does not say how rough is rough enough. Second,
does the theory need a determinacy threshold at all, a minimum net force below
which a case counts as a gap rather than as decided? On this Berman is
explicit that an aggregative system ``could incorporate thresholds or
eschew them,'' committing to neither.\footnote{Berman, \textit{Practices},
\textit{supra} note 7, at 331.} Any requirement of principled positivism
beyond the balance favoring one side is left open.

Third, is additive weighing the right functional form? On Berman's account,
principles combine by weight and magnitude rather than by fixed priority, a
structure he contrasts pointedly with the \textit{lexical} orderings that
rank some considerations absolutely above others.\footnote{\textit{Id.}\ at 311--13,
333 (distinguishing genuine weighing from lexical priority and locating
principled positivism on the weighing side of that line).} He does not defend the additive form against its nonlexical cousins,
though, nor does he say what turns on the contrast with counting or
ranking.

Fourth, and most consequentially, what \textit{drives changes} in the
practice that grounds the weights? On Berman's telling the weights are
grounded in practice,\footnote{\textit{Id.}\ at 329--30. The compressed
term ``practice-grounded'' is this Article's; Berman writes of principles
``grounded in practice.''} and he indicates, if
gesturally, how the taking-up and embedding of a principle fix its weight at
a given time. He leaves open the prior question of what moves those
takings-up: the shifting uptake of principles in judicial reasoning, their
embedding in holdings still relied upon, strategic adaptation by judges, the
turnover of the bench through appointment, or some combination.

Philosophical argument alone will not settle these, and the reason is that
each turns on a quantity. How much weight imprecision the model absorbs
depends on the relation among weight uncertainty, the number of principles,
and the force required for determinacy, across a parameter space too large
for any thought experiment to chart. The claim that the model yields more law than Hart can be
settled only by implementing both and measuring. Ruth Chang's objection
takes exactly that form, and it is the reason Part IV is built as a count
rather than as a worked case. To choose among
functional forms one must run the alternatives on identical inputs. And the
emergence of weights from practice can be tested only by simulating practice
and inspecting the resulting weights.

Computation enters at that point, to extend philosophical argument into
territory argument alone cannot reach. Given a formal model, a philosophical
theory can be stress-tested: implemented in code, subjected to systematic
variation, and measured against explicit criteria. The method will not
establish the truth of Berman's theory, a matter that would require
empirical validation well beyond this Article's scope. It shows instead how
far the theory's formal structure can bear the weight Berman places on
it.\footnote{Computational methods in jurisprudence remain uncommon. For
related uses of simulation in legal theory, see Frederick Schauer,
\textit{(Re)Taking Hart}, 119 \textsc{Harv.\ L.\ Rev.}\ 852, 879--82 (2006)
(discussing the value of formal modeling in positivist theorizing); Brian
Leiter, \textsc{Naturalizing Jurisprudence} 183--99 (2007) (arguing for
empirical methods in general jurisprudence).}

\subsection*{The Argument}

This Article develops the first systematic computational test of principled
positivism. It implements the theory faithfully, varies its parameters
across ranges no thought experiment could canvass, and measures the results
against what the theory itself requires. Two findings emerge.

The first is that the formal apparatus is sound. Given faithful weights,
meaning weights that correctly record how much each principle counts, the
weighted sum does what Berman says it does. Its sharpest result holds as
a theorem, before any weight is tuned: where a Hartian validation rule
contradicts itself, the weighted sum returns a single verdict. Such a rule
asks only for each authority's validity and direction, so on a genuine
conflict it certifies both sides at once. The sum cannot, because summing weighted forces yields one net force with one
sign.

Two further properties follow. The sum weighs the \textit{magnitude} of
competing considerations rather than ranking them by fixed priority, and it
tolerates rough weights, even wholly equal ones, without collapsing into
gridlock or into vote-counting. Several of these properties are theorems,
true of the model as such rather than artifacts of a particular
run.\footnote{See infra Part IV.D (establishing that the weighted sum is single-valued on any case whatever,
that the weight the embedding dynamic confers, which Part III implements as
the reliance mechanism, does not depend on how often a principle is invoked, and that equal weights do not reduce to counting
votes; each result is stated there with the scope it is proved under, which
differs across the five).}

How much law results remains a matter of degree, and there the
figures are conservative. At a demanding determinacy threshold the sum decides about six cases in
seven overall, and a slim majority of the hardest contested ones; on the
bare sign of net force, nearly every case but those in true equipoise. Both rates assume independent
activations, and correlation should only improve them. On the question Berman's critics press hardest, whether a balance of
principles can yield real law, the answer appears to be yes.

The second finding is the harder one. The formal structure does not fail. But
its success is inherited. Everything the model delivers rests on the
weights, and the weights rest on practice. How practice grounds a weight at
a given time the theory does say, if gesturally: officials take a principle
up and embed it in their reasoning and their holdings. It leaves open
the prior question of what \textit{drives changes} in those takings-up
over time.

Four such dynamics are modeled here: the rising and falling \textit{uptake}
of a principle in judicial reasoning, its \textit{embedding} in the holdings
that go on relying on it, the \textit{strategic adaptation} of judges
learning which arguments move their colleagues, and the near-random
\textit{shock} of judicial appointment. Each satisfies every requirement the
theory imposes. Each is well-defined, tolerant of rough weights,
magnitude-weighing, and more determinate than a hierarchy of the kind Hart's picture suggests, which
sorts the principles into ranked tiers and decides a case at the highest tier
holding an applicable principle, and only where enough of that tier's
principles point the same way.

But on the weights themselves the four do not converge. Applied to the same
cases, they disagree on the verdict in roughly a third and reach
opposite determinate verdicts in about a quarter.\footnote{Part VI reports
these divergence rates. The pattern has a telling exception: the appointment
dynamic produces a weight structure so close to embedding's that the two
disagree on fewer than one case in fifty, though that convergence turns out
to depend on both reading a common legal culture.} The choice among them is therefore unlikely to be a matter of legal
indifference. On this account it would shape the law's content.

Principled positivism, then, does not so much resolve the open question it
inherited as relocate it. Hart's law ran out in hard cases. Dworkin's
required moral interpretation to recover the right answer. Berman's locates
the answer in a balance of practice-grounded principles, and the balance
works. The theory does not settle which dynamics of changing grounds
fix the weights the sum runs on. Because those dynamics pull the law in different directions, the question is
no idle bookkeeping. It is itself a determinant of legal content, and one
that any successor theory would likely have to settle. That is the problem of practice, and it is likely a narrower and more
tractable problem than the one principled positivism inherited.

\subsection*{What the Argument Does and Does Not Establish}

Two cautions frame what follows. First, this is no refutation of principled
positivism. If anything the formal results vindicate it against its
sharpest critics, since the aggregation it proposes proves coherent, robust,
and more powerful than Hart's. And the claim here is narrower and, I think, more useful. The theory's reliance on practice is where its remaining work lies,
and that reliance is more consequential than the phrase
``practice-grounded'' lets on, a term this Article uses as shorthand for
his ``grounded in practice.'' Berman himself has pressed several of these
points, among them that the theory may not need a threshold at all, and that
the dynamics of change in legal grounds, including the near-random shock
of a new judicial appointment, are unlikely to be exhausted by any single
model. The analysis here takes those points seriously and shows,
computationally, why they matter.

Second, the test is formal rather than empirical. The cases are synthetic
and the practice is simulated. The model can show what the theory's structure will and will not
support, and it can perhaps suggest where a defender should look next. It cannot show whether any actual legal system grounds its weights through one
dynamic rather than another. Which practice best describes real adjudication is an empirical
question, a matter for the sociology of legal institutions and one that would
presumably require coding real judicial decisions, and this Article frames it
without resolving it.

\subsection*{Roadmap}

Part II sets out principled positivism: the SAW model, the roles of weight
and activation, and what the theory needs from each. Part III describes the
computational framework, the four dynamics of practice, and the four
requirements against which every result is measured, namely that the model
be well-defined, tolerate rough weights, weigh magnitude rather than rank by
fixed priority, and deliver more law than the tiered hierarchy that stands in
for Hart.

Part IV establishes the first finding. Given faithful weights the static
model works, and several of its properties are theorems rather than
artifacts of a run. Parts V and VI then turn to practice. Part V shows that
weights can \textit{emerge} from simulated practice rather than waiting on a
theorist to stipulate them. In Part VI the Article's central claim lands:
the four dynamics move the weights into distinct structures, and every one
of them meets the requirements the theory imposes. Practice, that is, is not
unitary.

Two robustness questions follow. Part VII shows that the result turns on
weighing magnitude rather than counting or ranking, and that within the
family of magnitude-weighing rules nothing turns on the particular form.
Part VIII takes up the determinacy threshold and shows it to be a modeling
choice rather than a commitment of the theory. It governs how much of the
law comes out determinate, not the model's substantive properties. Part IX then argues what the relocation would mean for the debate Berman
entered, and answers the objections it is likely to invite. Part X closes.

On their own, the computational findings reported here settle no question of
jurisprudence. What they do is sharpen one. The aggregation formula works,
and weights can emerge from practice rather than waiting on a theorist to
supply them, so the formal structure carries what Berman asks of it. That
structure cannot supply an account of the dynamics that move the
practice beneath it. Because those dynamics pull the weights in different
directions, the choice among them is itself a determinant of legal content.
That is the problem of practice: a narrower and, plausibly, a more
answerable problem than the one principled positivism inherited from
Hart and Dworkin, and arguably the right place for the next round of work to
begin.
